\section{Actual component sums in canonical principalization}

\begin{definition}[Finite sum of genuine connected boundary components]
\label{hir:finite-connected-component-sum-contract}
\uses{hir:connected-boundary-components,hir:component-product-effective-sum}
\lean{Hironaka.HironakaCompleteDivisorSum.HasFiniteComponentSum}
\leanok
For an effective divisor $D$ and a finite SNC boundary $E$, a finite component
sum consists of a distinct finite list $F$ of actual connected components of
members of $E$, still SNC, and a weighted list $A$ with underlying list $F$.
Each member is a closed smooth hypersurface and
\[
 D=\sum_{(C,a)\in A}a[C],\qquad
 \operatorname{divExp}(\mathcal I_D,C,x)=a\quad(x\in C).
\]
The sum is the actual sum of effective Cartier divisors constructed by
ideal products, and the equality is equality of those divisors. A coefficient
may be zero when a boundary component merely meets the support of $D$.
\end{definition}

\begin{theorem}[Complete sum on a finite connected SNC list]
\label{hir:complete-finite-connected-divisor-sum}
\uses{hir:component-factorization-original-coefficients,hir:component-product-effective-sum}
\lean{Hironaka.HironakaCompleteDivisorSum.exists_componentSum}
\leanok
Let $F$ be a distinct finite SNC list of closed connected smooth
hypersurfaces. Every effective divisor supported in $|F|$ is its actual
weighted component sum, with weights equal to the original divisor exponents
at every point of each component.
\end{theorem}
\begin{proof}
The complete factorization producer constructs the weighted list $A$, proves
$\mathcal I_D=\prod_{(C,a)\in A}\mathcal I_C^a$, and computes its coefficients
as the original exponents. The supplied actual closed codimension-one
geometry constructs each Cartier divisor $[C]$. The genuine finite sum has
exactly this ideal product. Equality of the invertible ideal determines the
effective divisor, giving the required sum equality and retaining the
computed exponents. Nothing about factorization or coefficient identification
is required from the caller.
\end{proof}

\begin{theorem}[Complete component sum for compact support]
\label{hir:compact-supported-complete-divisor-sum}
\uses{hir:complete-finite-connected-divisor-sum,hir:compact-connected-boundary-refinement}
\lean{Hironaka.HironakaCompleteDivisorSum.hasFiniteComponentSum_of_isCompact_support}
\leanok
Let $D$ have compact support contained in an arbitrary finite SNC boundary
$E$. Then $D$ has a finite component sum in the preceding sense, even if
the ambient manifold and the original boundary members are noncompact or
disconnected. This holds also in dimension zero.
\end{theorem}
\begin{proof}
In positive ambient dimension, the actual connected-component family of $E$
is locally finite. Compactness selects a finite distinct list of its members
meeting $\operatorname{Supp}D$. The coordinate refinement theorem makes that
list SNC and each member a closed connected smooth hypersurface. The union
of the selected list agrees with the original boundary on the compact
support, so it contains the support of $D$. Apply the preceding complete
sum producer. In dimension zero, no SNC member can contain a point: its
adapted singleton coordinate would lie in the empty coordinate index set.
Thus $\operatorname{Supp}D$ is empty, its actual ideal is the unit ideal,
and $D$ is the empty sum. Empty ambient spaces and empty lists are covered
by the same argument.
\end{proof}

\begin{theorem}[Canonical principalization with the actual divisor sum]
\label{hir:canonical-principalization-complete-divisor-sum}
\uses{hir:canonical-elementary-principalization-boundary,hir:compact-supported-complete-divisor-sum}
\source{wlodarczyk05-hironaka}
\lean{Hironaka.HironakaCompleteDivisorSum.exists_canonical_principalization_componentSum}
\leanok
Canonical principalization of a coherent ideal with nonzero stalks and
compact zero locus, relative to any initial SNC boundary, has its actual
effective divisor equal to a finite sum of genuine connected components
of the newly created boundary. Its coefficients are the original divisor
exponents. The same tower retains admissibility, the complete pullback ideal
identity, both support identities, the isomorphism off the zero locus,
canonicality, open restriction, and biholomorphism compatibility.
\end{theorem}
\begin{proof}
Construct the actual canonical principalization tower with all these
conclusions. Its projection is proper, so the preimage of the compact zero
locus is compact. This is exactly the support of its effective divisor.
The newly created boundary is the sublist obtained by dropping the
transforms of the initial members; a sublist of an SNC list is SNC, and
its union is the actual divisor support. Apply the compact-support sum
producer to that divisor and that sublist. Append its genuine component
sum to the conclusions of the very same tower. This uses no compactness
of the initial boundary and no assumption identifying product coefficients.
\end{proof}
